\chapter{Atributos de Clase y Métodos Especiales}

\section{Atributos de Clase vs Atributos de Instancia}
En Python existen dos tipos de atributos según su ámbito de pertenencia:

\begin{itemize}[leftmargin=*]
  \item \textit{Atributos de instancia}: pertenecen a un objeto específico y son inicializados habitualmente
    dentro del método \texttt{\_\_init\_\_} mediante \texttt{self.atributo}. Cada instancia mantiene su propia copia.
  \item \textit{Atributos de clase}: se definen directamente en el cuerpo de la clase (fuera de cualquier método)
    y son compartidos por todas las instancias creadas a partir de dicha clase.
\end{itemize}

\begin{lstlisting}[caption={Diferencia entre atributos de clase y de instancia.}, label={list:atributos_clase}]
class Student:
    # Atributo de clase (compartido por todos los estudiantes)
    school_name = "Instituto Técnico Salesiano Villada"

    def __init__(self, name, age):
        # Atributos de instancia (unicos para cada objeto)
        self.name = name
        self.age = age

e1 = Student("Lucas", 16)
e2 = Student("Sofia", 16)

print(e1.school_name)  # Muestra el valor de clase
print(e2.school_name)

# Modificacion mediante el nombre de la clase
Student.school_name = "ITSV Cordoba"
print(e1.school_name)  # Refleja el cambio en todas las instancias
\end{lstlisting}

Si se modifica un atributo de clase utilizando la referencia del nombre de la clase (\texttt{Student.school\_name}),
el cambio impacta en todas las instancias que no hayan sobrescrito dicho atributo a nivel de instancia.

\section{Métodos Especiales (Dunder Methods)}
Los métodos especiales, conocidos popularmente como \textit{dunder methods} (por la abreviatura de \textit{double underscore}),
son métodos cuyo nombre inicia y finaliza con dos guiones bajos, como por ejemplo \texttt{\_\_init\_\_}.
Python los utiliza para sobrecargar operadores y personalizar comportamientos integrados del lenguaje.

\subsection{Representación como Cadena: \texttt{\_\_str\_\_} y \texttt{\_\_repr\_\_}}
Cuando se imprime un objeto mediante \texttt{print()} o se convierte a cadena con \texttt{str()}, Python busca
el método \texttt{\_\_str\_\_}. Cuando se inspecciona en la consola interactiva o dentro de colecciones, se utiliza \texttt{\_\_repr\_\_}.

\begin{itemize}[leftmargin=*]
  \item \texttt{\_\_str\_\_}: pensado para presentar una salida amigable y legible para el usuario final.
  \item \texttt{\_\_repr\_\_}: pensado para proporcionar una representación inequívoca, útil para desarrolladores durante el proceso de depuración.
\end{itemize}

\begin{lstlisting}[caption={Implementación de \texttt{\_\_str\_\_} y \texttt{\_\_repr\_\_}.}, label={list:str_repr}]
class Book:
    def __init__(self, title, author):
        self.title = title
        self.author = author

    def __str__(self):
        return f"'{self.title}' escrito por {self.author}"

    def __repr__(self):
        return f"Book(title='{self.title}', author='{self.author}')"

libro = Book("Python Crash Course", "Eric Matthes")
print(str(libro))   # Invoca __str__
print(repr(libro))  # Invoca __repr__
\end{lstlisting}

\subsection{Sobrecarga de Operadores Aritméticos: \texttt{\_\_add\_\_}}
La sobrecarga de operadores permite definir cómo se comportarán los operadores matemáticos estándar
(como \texttt{+}, \texttt{-}, \texttt{*}) cuando se apliquen sobre objetos de nuestras propias clases.

\begin{lstlisting}[caption={Sobrecarga del operador suma mediante \texttt{\_\_add\_\_}.}, label={list:vector_add}]
class Vector2D:
    def __init__(self, x, y):
        self.x = x
        self.y = y

    def __add__(self, other):
        if isinstance(other, Vector2D):
            return Vector2D(self.x + other.x, self.y + other.y)
        return NotImplemented

    def __str__(self):
        return f"Vector({self.x}, {self.y})"

v1 = Vector2D(2, 4)
v2 = Vector2D(3, 1)
v3 = v1 + v2  # Invoca v1.__add__(v2)
print(v3)     # Salida: Vector(5, 5)
\end{lstlisting}

\subsection{Métodos Especiales \texttt{\_\_len\_\_} y \texttt{\_\_call\_\_}}
El método \texttt{\_\_len\_\_} permite a una clase responder a la función integrada \texttt{len()}.

\begin{lstlisting}[caption={Uso de \texttt{\_\_len\_\_} para calcular la cantidad de elementos.}, label={list:cart_len}]
class ShoppingCart:
    def __init__(self):
        self.items = []

    def add_item(self, item):
        self.items.append(item)

    def __len__(self):
        return len(self.items)

cart = ShoppingCart()
cart.add_item("Manzana")
cart.add_item("Pan")
print(f"Cantidad de productos: {len(cart)}")
\end{lstlisting}

Por otro lado, el método \texttt{\_\_call\_\_} permite que una instancia de una clase sea invocada como si fuera una función.

\begin{lstlisting}[caption={Instancia invocable con \texttt{\_\_call\_\_}.}, label={list:call_method}]
class Multiplier:
    def __init__(self, factor):
        self.factor = factor

    def __call__(self, x):
        return x * self.factor

duplicar = Multiplier(2)
print(duplicar(5))  # Retorna 10
\end{lstlisting}
